\documentclass[12pt]{article}
\usepackage{amsmath, amssymb, booktabs}

\title{Lecture 18: Marginal PMF, Correlation, Covariance, Joint Gaussian RVs}
\author{}
\date{}

\begin{document}

\maketitle

\section*{1. Marginal PMF}

\subsection*{1.1 Definition}

For discrete random variables $X$ and $Y$, the joint PMF is:
\[
p_{XY}(x,y) = P(X = x, Y = y)
\]

The marginal PMFs are obtained by summing over the other variable:
\[
p_X(x) = \sum_y p_{XY}(x,y)
\]
\[
p_Y(y) = \sum_x p_{XY}(x,y)
\]

\textbf{Key Idea}

The marginal PMF of one variable is obtained by summing the joint PMF over all possible values of the other variable.

\subsection*{1.2 Marginal PMF from Joint PMF Table}

From a joint PMF table:
\begin{itemize}
    \item Row sums $\rightarrow p_X(x)$
    \item Column sums $\rightarrow p_Y(y)$
\end{itemize}

Formally:
\[
p_X(x_i) = \sum_j p_{ij}, \quad p_Y(y_j) = \sum_i p_{ij}
\]

Also:
\[
\sum_i p_X(x_i) = 1, \quad \sum_j p_Y(y_j) = 1
\]

\subsection*{1.3 Example}

Given joint PMF:

\begin{center}
\begin{tabular}{c|cc}
 & $y=0$ & $y=1$ \\
\hline
$x=0$ & $1/8$ & $1/4$ \\
$x=1$ & $3/8$ & $1/4$ \\
\end{tabular}
\end{center}

\textbf{Step-by-step solution}

\textbf{Compute $p_X(x)$:}
\[
p_X(0) = \frac{1}{8} + \frac{1}{4} = \frac{3}{8}
\]
\[
p_X(1) = \frac{3}{8} + \frac{1}{4} = \frac{5}{8}
\]

\textbf{Compute $p_Y(y)$:}
\[
p_Y(0) = \frac{1}{8} + \frac{3}{8} = \frac{1}{2}
\]
\[
p_Y(1) = \frac{1}{4} + \frac{1}{4} = \frac{1}{2}
\]

\textbf{Final Answer}
\[
p_X(x) =
\begin{cases}
3/8 & x=0 \\
5/8 & x=1
\end{cases}
\quad
p_Y(y) =
\begin{cases}
1/2 & y=0 \\
1/2 & y=1
\end{cases}
\]

\section*{2. Correlation}

\subsection*{2.1 Definition}

\[
R_{XY} = E[XY]
\]

It represents the average value of the product $XY$.

\subsection*{2.2 Interpretation and Limitation}

\begin{itemize}
    \item $R_{XY}$ depends on the scale of $X$ and $Y$
    \item Computed about the origin, not about means
    \item Mixes mean values, spread, and dependence
\end{itemize}

\textbf{Important Result}
\[
R_{XY} = 0 \Rightarrow \text{orthogonal about the origin}
\]

\textbf{Transition}

To remove mean effects $\rightarrow$ center variables $\rightarrow$ covariance

\section*{3. Covariance}

\subsection*{3.1 Definition}

\[
\text{Cov}(X,Y) = E[(X - \mu_X)(Y - \mu_Y)]
\]

\[
\mu_X = E[X], \quad \mu_Y = E[Y]
\]

Measures how centered variables vary together.

\subsection*{3.2 Interpretation Summary}

\begin{itemize}
    \item $\text{Cov}(X,Y) > 0$ : positive covariance
    \item $\text{Cov}(X,Y) < 0$ : negative covariance
    \item $\text{Cov}(X,Y) = 0$ : no linear co-variation
\end{itemize}

Indicates direction only; magnitude depends on units.

\subsection*{3.3 Algebraic Form (Derivation)}

\[
\text{Cov}(X,Y) = E[(X - \mu_X)(Y - \mu_Y)]
\]

\textbf{Step 1: Expand}
\[
= E[XY - X\mu_Y - \mu_X Y + \mu_X \mu_Y]
\]

\textbf{Step 2: Apply linearity}
\[
= E[XY] - \mu_Y E[X] - \mu_X E[Y] + \mu_X \mu_Y
\]

\textbf{Step 3: Substitute means}
\[
= E[XY] - \mu_X \mu_Y
\]

\textbf{Final Result}
\[
\boxed{\text{Cov}(X,Y) = E[XY] - E[X]E[Y]}
\]

\subsection*{3.4 Properties}

\begin{enumerate}
    \item $\text{Cov}(X,Y) = \text{Cov}(Y,X)$
    \item $\text{Cov}(X,X) = \text{Var}(X)$
    \item $\text{Cov}(aX + b, cY + d) = ac \cdot \text{Cov}(X,Y)$
    \item $X \perp Y \Rightarrow \text{Cov}(X,Y) = 0$
\end{enumerate}

Zero covariance does NOT imply independence.

\subsection*{3.5 Example 1}

Given:
\[
Y = -2X + 3
\]

\[
\text{Cov}(X,Y) = E[XY] - E[X]E[Y]
\]

\[
XY = -2X^2 + 3X
\]

\[
E[XY] = -2E[X^2] + 3E[X]
\]

\[
E[Y] = -2E[X] + 3
\]

\[
\text{Cov}(X,Y) = -2E[X^2] + 2(E[X])^2
\]

\[
\text{Var}(X) = E[X^2] - (E[X])^2
\]

\[
\Rightarrow \text{Cov}(X,Y) = -2\text{Var}(X)
\]

\section*{4. Pearson’s Correlation Coefficient}

\subsection*{4.1 Definition}

\[
\rho_{XY} = \frac{\text{Cov}(X,Y)}{\sigma_X \sigma_Y}
\]

\subsection*{4.2 Properties}

\[
-1 \leq \rho_{XY} \leq 1
\]

Measures strength and direction of linear relationship.

\subsection*{4.3 Interpretation}

\begin{itemize}
    \item $\rho_{XY} > 0$ : positive trend
    \item $\rho_{XY} < 0$ : negative trend
    \item $\rho_{XY} = 0$ : no linear correlation
\end{itemize}

\section*{5. Example 2 (Continuous Case)}

\[
f_{XY}(x,y) =
\begin{cases}
\frac{xy}{96}, & 0<x<4,\; 1<y<5 \\
0, & \text{otherwise}
\end{cases}
\]

\[
E[X] = \frac{8}{3}, \quad E[Y] = \frac{31}{9}
\]

\[
E[XY] = \frac{248}{27}
\]

\[
E[2X+3Y] = \frac{47}{3}
\]

\[
\text{Var}(X) = \frac{8}{9}, \quad \text{Var}(Y) = \frac{92}{81}
\]

\[
\text{Cov}(X,Y) = 0, \quad \rho_{XY} = 0
\]

\section*{6. Jointly Gaussian Random Variables}

\subsection*{6.1 Definition}

A pair $(X,Y)$ is jointly Gaussian if:
\begin{itemize}
    \item Joint density is bivariate Gaussian
    \item Marginals are Gaussian
\end{itemize}

\subsection*{6.3 Joint Gaussian PDF}

\[
f_{XY}(x,y) =
\frac{1}{2\pi \sigma_X \sigma_Y \sqrt{1-\rho^2}}
\exp\left(
-\frac{1}{2(1-\rho^2)}
\left[
\left(\frac{x-\mu_X}{\sigma_X}\right)^2 +
\left(\frac{y-\mu_Y}{\sigma_Y}\right)^2 -
2\rho \left(\frac{x-\mu_X}{\sigma_X}\right)\left(\frac{y-\mu_Y}{\sigma_Y}\right)
\right]
\right)
\]

\section*{Final Key Takeaways}

\begin{enumerate}
    \item Marginal PMF = sum over joint PMF
    \item Correlation is raw and not centered
    \item Covariance removes mean effect
    \item $\text{Cov}(X,Y) = E[XY] - E[X]E[Y]$
    \item Zero covariance $\neq$ independence
    \item Pearson coefficient normalizes covariance
    \item Joint Gaussian fully described by means, variances, and correlation
\end{enumerate}

\end{document}